% ══════════════════════════════════════════════════════════════
% Chapter 9 — Feature Engineering from Risk Systems
% ══════════════════════════════════════════════════════════════

\chapter{Feature Engineering from Risk Systems}
\label{ch:features}

\begin{projectconnection}[Project Connection --- Phase 0 \& Phase 2]
This chapter maps directly to Phase~0 (data pipeline) and Phase~2
(feature engineering).  Every feature family here corresponds to a
module in \texttt{src/features/}.  The critical lesson: \textbf{point-in-time
discipline}.  A single lookahead bug---using VaR from the overnight risk
run before it was actually available---invalidates your entire backtest.
\end{projectconnection}


% ──────────────────────────────────────────────────────────────
\section{Point-in-Time Discipline}
\label{sec:pit}
% ──────────────────────────────────────────────────────────────

\begin{prereq}[Feature Matrix Notation]
A feature matrix $\bX \in \R^{T \times p}$ has $T$ rows (one per time
step) and $p$ columns (one per feature).  Row $t$ contains only
information that was \emph{known and available} at time $t$.  If any
entry $X_{t,j}$ uses data that arrived after $t$, the matrix is
contaminated by lookahead bias.
\end{prereq}

Every feature must be stamped with \emph{when it was known}, not when
it applied.  Lopez de Prado (2018, Ch.~2) calls this
\textbf{point-in-time} (PIT) stamping and identifies lookahead bias as
the single most common and most damaging error in financial ML.
The principle is simple: a feature available only at time $t+\delta$
cannot appear in any row earlier than $t+\delta$ in your feature
matrix.

\subsection{Knowledge-Date Lagging}

SecDB runs its nightly risk engine in a batch process.  The VaR number
for business date $T$ is computed overnight between $T$ close and $T+1$
open.  Therefore:

\begin{itemize}[nosep]
  \item The \textbf{as-of date} of VaR is $T$---it measures the
        portfolio's risk as of the close on date $T$.
  \item The \textbf{knowledge date} is $T+1$ morning---this is when the
        number becomes available to traders and models.
  \item Any model using VaR$_T$ as a feature may only do so for
        predictions dated $T+1$ or later.
\end{itemize}

The gap between as-of date and knowledge date varies by data source:

\begin{center}
\small
\begin{tabular}{@{}lll@{}}
\toprule
\textbf{Risk Output} & \textbf{As-of Date} & \textbf{Knowledge Date} \\
\midrule
Nightly VaR / ES          & $T$ close  & $T+1$ pre-open \\
Scenario P\&L              & $T$ close  & $T+1$ pre-open \\
Intraday Greeks            & Real-time  & Real-time \\
VaR limit (set by risk)    & Set date   & Set date \\
Factor-VaR decomposition   & $T$ close  & $T+1$ pre-open \\
\bottomrule
\end{tabular}
\end{center}

\noindent Intraday Greeks are the exception: they are computed in
real-time by SecDB's pricing engine and available immediately.  This is
why Project~2 (Chapter~\ref{ch:microstructure}, Phase~4B) can use
same-day Greek features.

\begin{definition}[Point-in-Time Stamping]
For a risk output $x$ with as-of date $T$ and knowledge date $T+\delta$,
the PIT-stamped feature is:
\begin{equation}
\label{eq:pit}
\tilde{x}_t =
\begin{cases}
x_T & \text{if } t \ge T + \delta \\
\texttt{NaN} & \text{if } t < T + \delta
\end{cases}
\end{equation}
where $\delta$ is the \textbf{publication lag}---typically 1 business day
for overnight batch outputs.  The tilde notation $\tilde{x}$ denotes the
PIT-lagged version.
\end{definition}

\begin{workedexample}{VaR Point-in-Time Timeline}
Consider the following sequence for a Monday--Wednesday window:

\begin{center}
\small
\begin{tabular}{@{}llll@{}}
\toprule
\textbf{Event} & \textbf{Wall-Clock Time} & \textbf{As-of} & \textbf{Knowledge} \\
\midrule
Monday close       & Mon 4:00pm  & ---     & --- \\
Risk batch starts  & Mon 6:00pm  & Monday  & --- \\
Risk batch ends    & Tue 5:00am  & Monday  & Tuesday \\
Tuesday open       & Tue 9:30am  & ---     & --- \\
Model predicts     & Tue 9:30am  & ---     & Uses VaR$_{\text{Mon}}$ \\
Tuesday close      & Tue 4:00pm  & ---     & --- \\
Risk batch ends    & Wed 5:00am  & Tuesday & Wednesday \\
\bottomrule
\end{tabular}
\end{center}

The model running at Tuesday open may use Monday's VaR (batch completed
at 5am Tuesday) but \emph{not} Tuesday's VaR (batch has not run yet).
In code, this means:

\begin{equation}
\label{eq:pit-example}
\tilde{\VaR}_{t=\text{Tue}} = \VaR_{\text{Mon}}, \qquad
\tilde{\VaR}_{t=\text{Wed}} = \VaR_{\text{Tue}}
\end{equation}

Using $\VaR_{\text{Tue}}$ at Tuesday close is cheating---it was not
available until Wednesday morning.
\end{workedexample}

\begin{warning}[A Single Lookahead Bug Invalidates Your Entire Backtest]
VaR for date $T$, computed in the overnight batch, is available at
$T+1$ open.  Using it at $T$ close is cheating.  The effect is subtle:
your backtest will show inflated Sharpe ratios because the model
``knows'' tomorrow's risk state today.  Lopez de Prado (2018, Ch.~2)
emphasizes that lookahead bias is the most frequent source of false
discoveries in financial ML---more common than overfitting and harder
to detect because the backtest looks clean.  In your pipeline code
(\texttt{src/data/point\_in\_time.py}), enforce PIT lagging at the data
ingestion layer so no downstream code can accidentally skip it.
\end{warning}

\subsection{Implementation Discipline}

Three rules for your \texttt{src/data/pipeline.py}:

\begin{enumerate}[nosep]
  \item \textbf{Lag at ingestion, not at modeling.}  Apply the
        knowledge-date shift when data enters the pipeline.  Every
        downstream consumer sees only PIT-valid data.
  \item \textbf{Store both timestamps.}  Keep the as-of date and
        knowledge date as separate columns.  This allows auditing.
  \item \textbf{Unit test with synthetic data.}  Create a test where a
        feature's knowledge date is deliberately set to $T+1$.  Verify
        that the feature at time $T$ is \texttt{NaN}, not the $T$ value.
\end{enumerate}


% ──────────────────────────────────────────────────────────────
\section{Feature Family 1: VaR Utilization}
\label{sec:feat-var-util}
% ──────────────────────────────────────────────────────────────

\begin{prereq}[VaR Utilization Recap]
VaR utilization $= \VaR_{\text{current}} / \VaR_{\text{limit}} \times
100\%$ measures how much of the desk's risk budget is deployed.
Desks typically operate at 50--80\%; approaching 100\% forces position
reduction.  See Section~\ref{sec:var-utilization} of
Chapter~\ref{ch:risk-systems} for the full definition and the
procyclical leverage mechanism (Adrian and Shin, 2014).
\end{prereq}

VaR utilization is one of two \textbf{priority} feature families
(design spec, Phase~2).  The theoretical basis is strong: VaR
utilization is the most direct proxy for the binding constraint in
intermediary asset pricing models (He and Krishnamurthy, 2013;
Chapter~\ref{ch:intermediary}).

\subsection{Feature Definitions}

Let $U_t = \VaR_t / \VaR_{\text{limit},t}$ denote VaR utilization at
time $t$ (PIT-lagged per Section~\ref{sec:pit}).  Three features:

\begin{definition}[VaR Utilization Features]
\textbf{(1) Level:}
\begin{equation}
\label{eq:util-level}
f^{\text{util-level}}_t = \tilde{U}_t
\end{equation}
Current utilization percentage.  Values above 85\% signal constrained
dealers.

\textbf{(2) Rate of change} over horizon $h \in \{1, 5, 21\}$ days:
\begin{equation}
\label{eq:util-roc}
f^{\text{util-}\Delta h}_t = \tilde{U}_t - \tilde{U}_{t-h}
\end{equation}
Rising utilization signals building stress; the 5-day change captures
weekly dynamics; the 21-day change captures monthly trends.

\textbf{(3) Z-score} over a rolling window of length $w$ (default $w=63$
business days, approximately 3 months):
\begin{equation}
\label{eq:util-zscore}
f^{\text{util-z}}_t = \frac{\tilde{U}_t - \bar{U}_{t,w}}
  {\hat{\sigma}_{U,t,w}}
\end{equation}
where $\bar{U}_{t,w}$ and $\hat{\sigma}_{U,t,w}$ are the rolling mean
and standard deviation of $\tilde{U}$ over the window
$[t-w+1,\; t]$.

\begin{itemize}[nosep]
  \item $f^{\text{util-level}}$ --- raw utilization; captures the
        constraint \emph{level}.
  \item $f^{\text{util-}\Delta h}$ --- first difference; captures the
        \emph{direction} of constraint tightening or loosening.
  \item $f^{\text{util-z}}$ --- standardized anomaly; captures how
        \emph{unusual} today's utilization is relative to recent history.
\end{itemize}
\end{definition}

\begin{workedexample}{VaR Utilization Features from a 10-Day Series}
Given PIT-lagged utilization data:

\begin{center}
\small
\begin{tabular}{@{}ccccc@{}}
\toprule
$t$ & $\tilde{U}_t$ (\%) & $\Delta 1$d & $\Delta 5$d &
  Z-score ($w\!=\!5$) \\
\midrule
1  & 62 & ---  & ---  & --- \\
2  & 64 & $+2$ & ---  & --- \\
3  & 68 & $+4$ & ---  & --- \\
4  & 71 & $+3$ & ---  & --- \\
5  & 73 & $+2$ & $+11$& --- \\
6  & 78 & $+5$ & $+14$& $+1.63$ \\
7  & 76 & $-2$ & $+12$& $+0.99$ \\
8  & 80 & $+4$ & $+12$& $+1.38$ \\
9  & 83 & $+3$ & $+12$& $+1.55$ \\
10 & 88 & $+5$ & $+15$& $+1.82$ \\
\bottomrule
\end{tabular}
\end{center}

\begin{itemize}[nosep]
  \item \textbf{Level:} At $t\!=\!10$, utilization is 88\%---above the
        85\% warning threshold.
  \item \textbf{$\Delta 5$d:} The 5-day change of $+15$ percentage
        points signals rapid constraint tightening.
  \item \textbf{Z-score:} At $+1.82$, today's utilization is nearly
        2 standard deviations above its recent 5-day rolling mean.
\end{itemize}
In practice, use $w=63$ (3 months) for the z-score to have a
meaningful rolling baseline.  The $w=5$ window here is for illustration
only.
\end{workedexample}

\begin{intuition}[Why Three Features, Not Just One?]
The level, rate of change, and z-score capture fundamentally different
information.  A desk at 85\% utilization that has been there for 6
months (z-score near 0) is in a different state than a desk that jumped
from 60\% to 85\% in a week (z-score near $+2$).  Both have the same
level, but the second desk is under acute stress and more likely to
be forced to sell.  The z-score distinguishes chronic from acute
constraint tightening.
\end{intuition}


% ──────────────────────────────────────────────────────────────
\section{Feature Family 2: Factor Concentration}
\label{sec:feat-factor-conc}
% ──────────────────────────────────────────────────────────────

\begin{prereq}[Factor-VaR Decomposition and HHI]
Factor-VaR decomposes total VaR into contributions from underlying risk
factors (rates, credit, equity, FX, commodities).  The
Herfindahl-Hirschman Index $\HHI = \sum_k s_k^2$ measures concentration,
where $s_k$ is factor $k$'s share of total VaR.  $\HHI \in [1/K, 1]$;
higher values mean more concentrated risk.  See
Section~\ref{sec:factor-var} and specifically equation~\eqref{eq:hhi}
of Chapter~\ref{ch:risk-systems}.
\end{prereq}

Factor concentration is the second \textbf{priority} feature family.
The theory: when many dealers crowd into the same factor, HHI rises.
If that factor reverses, concentrated dealers must all deleverage
simultaneously---triggering fire sales (Coval and Stafford, 2007)
and predictable return reversals (He, Kelly, and Manela, 2017;
Chapter~\ref{ch:intermediary}).

\subsection{Feature Definitions}

Let $s_{k,t}$ denote the PIT-lagged factor-VaR share of factor $k$ at
time $t$.

\begin{definition}[Factor Concentration Features]
\textbf{(1) Herfindahl Index:}
\begin{equation}
\label{eq:feat-hhi}
f^{\text{HHI}}_t = \sum_{k=1}^{K} \tilde{s}_{k,t}^{\,2}
\end{equation}
Measures overall concentration.  $K$ is the number of risk factors
(typically 5--10 for a cross-asset desk).

\textbf{(2) Top-3 factor share:}
\begin{equation}
\label{eq:top3}
f^{\text{top3}}_t = \sum_{k \in \mathcal{T}_t} \tilde{s}_{k,t}
\end{equation}
where $\mathcal{T}_t$ is the set of three factors with the largest
VaR shares at time $t$.  This is a simpler, more interpretable
concentration metric.

\textbf{(3) Rolling HHI change:}
\begin{equation}
\label{eq:delta-hhi}
f^{\Delta\text{HHI}}_t = f^{\text{HHI}}_t - f^{\text{HHI}}_{t-h}
\end{equation}
for $h \in \{5, 21\}$ days.  Rising HHI means concentration is
increasing.

\begin{itemize}[nosep]
  \item $f^{\text{HHI}}$ --- level of concentration across all factors.
  \item $f^{\text{top3}}$ --- concentration in the dominant factors;
        robust to noise in small factor shares.
  \item $f^{\Delta\text{HHI}}$ --- direction and speed of
        concentration changes.
\end{itemize}
\end{definition}

\begin{keyidea}[High HHI = Concentrated Risk = Fragility]
High HHI means one factor dominates the risk profile.  This signals
crowding and fragility.  He, Kelly, and Manela (2017) show that a single
intermediary-capital factor prices risk across asset classes---when
dealers are concentrated in one factor, the cross-asset pricing kernel
is exposed to a single point of failure.

\textbf{Concrete numbers:} For a 5-factor desk, $\HHI = 0.20$ means
all factors contribute equally ($N_{\text{eff}} = 5$).  $\HHI = 0.50$
means the desk is effectively a 2-factor bet.  $\HHI > 0.60$ should
raise concern about crowding.
\end{keyidea}


% ──────────────────────────────────────────────────────────────
\section{Feature Family 3: VaR Dynamics}
\label{sec:feat-var-dynamics}
% ──────────────────────────────────────────────────────────────

\begin{prereq}[VaR Computation Methods and the Cliff Effect]
VaR under historical simulation suffers from the cliff effect: an
extreme observation entering or leaving the lookback window can cause
a discontinuous jump (Section~\ref{subsec:hist-sim} of
Chapter~\ref{ch:risk-systems}).  VaR-dynamics features should be
smoothed (e.g., 5-day moving average) to reduce this noise.
\end{prereq}

VaR dynamics capture the \emph{time-series behavior} of the aggregate
risk number---not where risk lives (that is factor concentration) or
how close the desk is to its limit (that is utilization), but how risk
itself is evolving.

\subsection{Feature Definitions}

Let $V_t$ denote the desk's total VaR at time $t$, PIT-lagged.

\begin{definition}[VaR Dynamics Features]
\textbf{(1) Delta VaR:}
\begin{equation}
\label{eq:delta-var}
f^{\Delta\text{VaR}}_t = \tilde{V}_t - \tilde{V}_{t-1}
\end{equation}
The daily change in VaR.  Large positive values indicate a rapid
increase in measured risk.

\textbf{(2) Component VaR by asset class:}
\begin{equation}
\label{eq:cvar-ac}
f^{\text{CVaR-}c}_t = \widetilde{\text{CVaR}}_{c,t}
\end{equation}
for each asset class $c \in \{\text{rates, credit, equity, FX,
commodities}\}$.  These track the risk contribution of each asset class
over time (see equation~\eqref{eq:cvar} of Chapter~\ref{ch:risk-systems}
for the component VaR definition).

\textbf{(3) VaR momentum:}
\begin{equation}
\label{eq:var-mom}
f^{\text{VaR-mom}}_t = \operatorname{sign}\!\Bigl(
  \sum_{i=0}^{4} (\tilde{V}_{t-i} - \tilde{V}_{t-i-1})
\Bigr) \cdot
\biggl|\sum_{i=0}^{4} (\tilde{V}_{t-i} - \tilde{V}_{t-i-1})\biggr|^{1/2}
\end{equation}
Sign and square-root magnitude of the 5-day rolling VaR change.  The
square root compresses extreme values while preserving sign.

\begin{itemize}[nosep]
  \item $f^{\Delta\text{VaR}}$ --- instantaneous risk change.
  \item $f^{\text{CVaR-}c}$ --- risk attribution by asset class.
  \item $f^{\text{VaR-mom}}$ --- direction and persistence of risk
        trends.
\end{itemize}
\end{definition}

\begin{warning}[Methodology Breaks Masquerade as Signals]
Section~\ref{sec:methodology-changes} of Chapter~\ref{ch:risk-systems}
warns that VaR model updates create structural breaks.  A switch from
historical simulation to Monte Carlo can produce a one-day $\Delta$VaR
that dwarfs any genuine portfolio change.  LightGBM
(Chapter~\ref{ch:trees}) will happily learn to split on these artifacts.
Before engineering VaR-dynamics features, obtain the methodology change
dates from Phase~0's risk-model audit and either include them as
binary control variables or restrict training to a single methodology
regime.
\end{warning}


% ──────────────────────────────────────────────────────────────
\section{Feature Family 4: Scenario P\&L}
\label{sec:feat-scenario}
% ──────────────────────────────────────────────────────────────

\begin{prereq}[Scenario Analysis Recap]
A scenario specifies simultaneous moves in multiple risk factors
(e.g., rates $+200$bp, equities $-25\%$).  The risk system revalues the
portfolio under each scenario to produce a scenario P\&L.  See
Section~\ref{sec:scenarios} of Chapter~\ref{ch:risk-systems} for
types of scenarios and a worked numerical example.
\end{prereq}

Scenario P\&L features are forward-looking in a way that VaR is not.
VaR uses backward-looking distributions; scenarios ask ``what would
happen \emph{if}?''  Changes in scenario P\&L capture shifts in
portfolio vulnerability that may not yet appear in VaR.

\subsection{Feature Definitions}

Let $\pi_{j,t}$ denote the PIT-lagged scenario P\&L for scenario $j$
at time $t$, with $j = 1, \ldots, J$ scenarios.

\begin{definition}[Scenario P\&L Features]
\textbf{(1) Worst-case rank:}
\begin{equation}
\label{eq:worst-rank}
f^{\text{worst-id}}_t = \arg\min_j \;\tilde{\pi}_{j,t}
\end{equation}
Which scenario produces the largest loss today?  This is a categorical
feature---the \emph{identity} of the worst scenario, not its magnitude.

\textbf{(2) Scenario dispersion:}
\begin{equation}
\label{eq:scenario-disp}
f^{\text{disp}}_t = \sqrt{\frac{1}{J}\sum_{j=1}^{J}
  \bigl(\tilde{\pi}_{j,t} - \bar{\pi}_t\bigr)^2}
\end{equation}
where $\bar{\pi}_t = J^{-1}\sum_j \tilde{\pi}_{j,t}$ is the
cross-scenario mean.  High dispersion means the portfolio is
directionally sensitive---performance varies sharply with the type of
stress.  Low dispersion means the portfolio is well-hedged or
indifferent to the specific scenario.

\textbf{(3) Worst-case identity change:}
\begin{equation}
\label{eq:worst-flip}
f^{\text{flip}}_t = \mathbf{1}\!\bigl[
  f^{\text{worst-id}}_t \ne f^{\text{worst-id}}_{t-1}\bigr]
\end{equation}
A binary indicator: did the worst scenario change from yesterday?
A flip signals a qualitative shift in portfolio vulnerability---the
desk's primary risk has rotated.

\textbf{(4) Scenario P\&L skewness:}
\begin{equation}
\label{eq:scenario-skew}
f^{\text{skew}}_t = \frac{J^{-1}\sum_{j=1}^{J}
  (\tilde{\pi}_{j,t} - \bar{\pi}_t)^3}
  {\bigl(f^{\text{disp}}_t\bigr)^3}
\end{equation}
Negative skewness means a few scenarios produce very large losses
relative to the rest---the portfolio has a fat left tail across
scenarios.

\begin{itemize}[nosep]
  \item $f^{\text{worst-id}}$ --- categorical; use one-hot encoding or
        target encoding for tree models.
  \item $f^{\text{disp}}$ --- continuous; high values signal directional
        exposure.
  \item $f^{\text{flip}}$ --- binary; captures qualitative regime shifts.
  \item $f^{\text{skew}}$ --- continuous; captures asymmetry in
        scenario outcomes.
\end{itemize}
\end{definition}


% ──────────────────────────────────────────────────────────────
\section{Feature Family 5: Cross-Asset Flow}
\label{sec:feat-cross-asset}
% ──────────────────────────────────────────────────────────────

\begin{prereq}[Component VaR Shares]
Component VaR (Section~\ref{sec:component-var} of
Chapter~\ref{ch:risk-systems}) decomposes total VaR by position or
asset class.  The component VaR share of asset class $c$ is
$s_{c,t} = \text{CVaR}_{c,t} / \VaR_t$.  Shares sum to 1 (Euler
decomposition).
\end{prereq}

Cross-asset flow features track how risk \emph{migrates} between asset
classes.  These features are distinct from factor concentration
(Section~\ref{sec:feat-factor-conc}), which measures overall
diversification.  Cross-asset flow measures the \emph{direction} of
risk rotation.

\subsection{Feature Definitions}

Let $s_{c,t}$ denote the PIT-lagged component VaR share of asset class
$c$ at time $t$.

\begin{definition}[Cross-Asset Flow Features]
\textbf{(1) Component VaR share shift:}
\begin{equation}
\label{eq:share-shift}
f^{\Delta s_c}_t = \tilde{s}_{c,t} - \tilde{s}_{c,t-h}
\end{equation}
for each asset class $c$ and horizon $h \in \{5, 21\}$ days.  Positive
values mean the desk is accumulating risk in asset class $c$ relative
to other classes.

\textbf{(2) Rolling cross-asset VaR correlation:}
\begin{equation}
\label{eq:cross-corr}
f^{\rho_{c_1,c_2}}_t = \operatorname{Corr}\!\bigl(
  \tilde{s}_{c_1,\, t-w+1:t},\; \tilde{s}_{c_2,\, t-w+1:t}
\bigr)
\end{equation}
over a rolling window of $w=63$ days, for key asset-class pairs (e.g.,
rates--credit, equity--commodities).  High positive correlation means
the desk's risk in both classes moves together; high negative
correlation suggests active rotation from one class to another.

\begin{itemize}[nosep]
  \item $f^{\Delta s_c}$ --- direction of risk rotation by asset class.
  \item $f^{\rho_{c_1,c_2}}$ --- persistence and co-movement of risk
        allocation decisions.
\end{itemize}
\end{definition}

\begin{keyidea}[Risk Rotation Predicts Relative Performance]
If the rates VaR share rises while the credit VaR share falls,
dealers are rotating risk---accumulating rate exposure and shedding
credit exposure.  Under the intermediary pricing framework
(Chapter~\ref{ch:intermediary}), this shift in dealer balance sheets
should predict \emph{relative} performance between rates and credit.
He, Kelly, and Manela (2017) show that a single intermediary-capital
factor prices the cross-section of returns across asset classes.
Cross-asset flow features test whether \emph{changes} in risk
allocation across classes carry incremental information beyond
the aggregate constraint level.
\end{keyidea}


% ──────────────────────────────────────────────────────────────
\section{Feature Family 6: Dealer Greeks (Project 2)}
\label{sec:feat-greeks}
% ──────────────────────────────────────────────────────────────

\begin{prereq}[Dealer Greeks]
A dealer's net gamma, vega, vanna, and charm determine whether hedging
flows stabilize or destabilize prices.  Short-gamma dealers amplify
price moves; long-gamma dealers dampen them.  See
Sections~\ref{sec:gamma} and~\ref{sec:higher-greeks} of
Chapter~\ref{ch:microstructure} for definitions and the
shock-absorber vs.\ shock-amplifier intuition.
\end{prereq}

These features belong to Phase~4B (Project~2: Book-Gamma Intraday
Momentum).  They are included here for completeness and because the
shared infrastructure (Task~5, data pipeline) should accommodate them
from the start.  All Greek features are available in real time from
SecDB's pricing engine, so PIT lagging is not required---but you must
still ensure that features used for end-of-day predictions are computed
at a consistent timestamp (e.g., 3:00pm snapshot).

\subsection{Feature Definitions}

Let $\Gamma^{\text{net}}_t$, $\mathcal{V}^{\text{net}}_t$,
$\text{Vanna}^{\text{net}}_t$, and $\text{Charm}^{\text{net}}_t$ denote
the aggregate net book-level Greeks from SecDB at time $t$.

\begin{definition}[Dealer Greek Features]
\textbf{(1) Net gamma, vega, vanna, charm:}
\begin{equation}
\label{eq:net-greeks}
f^{\Gamma}_t = \Gamma^{\text{net}}_t, \quad
f^{\mathcal{V}}_t = \mathcal{V}^{\text{net}}_t, \quad
f^{\text{vanna}}_t = \text{Vanna}^{\text{net}}_t, \quad
f^{\text{charm}}_t = \text{Charm}^{\text{net}}_t
\end{equation}
These are the raw aggregate sensitivities.  The sign of gamma is the
key feature: negative aggregate gamma predicts intraday momentum
(Baltussen, Da, Lammers, and Martens, 2021;
Section~\ref{sec:intraday-momentum} of Chapter~\ref{ch:microstructure}).

\textbf{(2) Gamma-flip level:}
\begin{equation}
\label{eq:gamma-flip}
S^{\text{flip}}_t = \text{strike at which aggregate } \Gamma^{\text{net}}
  \text{ changes sign}
\end{equation}
The gamma-flip level is the underlying price where the desk transitions
from long gamma (stabilizing) to short gamma (destabilizing) or vice
versa.

\textbf{(3) Distance to flip:}
\begin{equation}
\label{eq:dist-flip}
f^{\text{d-flip}}_t = \frac{S_t - S^{\text{flip}}_t}{S_t} \times 100\%
\end{equation}
where $S_t$ is the current underlying price.  Small absolute values
mean the market is near the flip level---a regime transition in hedging
flow direction is imminent.

\begin{itemize}[nosep]
  \item $f^{\Gamma}, f^{\mathcal{V}}, f^{\text{vanna}},
        f^{\text{charm}}$ --- raw book-level sensitivities with
        correct dealer sign (not estimated from public data).
  \item $S^{\text{flip}}$ --- the gamma regime boundary.
  \item $f^{\text{d-flip}}$ --- proximity to regime transition;
        nonlinear effects concentrate near zero.
\end{itemize}
\end{definition}


% ──────────────────────────────────────────────────────────────
\section{Confound Checking}
\label{sec:confounds}
% ──────────────────────────────────────────────────────────────

\begin{prereq}[Omitted Variable Bias]
If a predictor $x$ and outcome $y$ are both driven by an omitted
variable $z$, the apparent predictive relationship between $x$ and $y$
is spurious.  Including $z$ as a control variable reveals whether $x$
has \emph{incremental} predictive power.  This is the financial-ML
analogue of the ``confound'' in causal inference.
\end{prereq}

A risk-system feature that merely repackages information already
available in public data adds no value.  VaR rises when markets are
stressed---but so does the VIX.  If your VaR-based signal vanishes
after controlling for VIX, you have not discovered a new signal; you
have built a noisy VIX proxy.

\subsection{Control Variables}

Add the following public factors as controls in every model:

\begin{enumerate}[nosep]
  \item \textbf{VIX level and change:} $\text{VIX}_t$ and
        $\Delta\text{VIX}_t$.  The most common confound for any
        risk-based feature.
  \item \textbf{Credit spread:} Investment-grade and high-yield OAS
        (option-adjusted spread).  Confounds credit-related features.
  \item \textbf{Term slope:} 10y$-$2y Treasury yield spread.  Confounds
        rate-related features and captures macro regime.
  \item \textbf{Equity market return:} S\&P 500 trailing 5-day return.
        Confounds any feature correlated with market direction.
  \item \textbf{Dollar index:} DXY level and change.  Confounds
        FX-related features.
\end{enumerate}

\subsection{The Confound Test}

Run two models (ridge baseline, then LightGBM---see
Chapters~\ref{ch:regularized} and~\ref{ch:trees}):

\begin{enumerate}[nosep]
  \item \textbf{Model A:} Risk features only.
  \item \textbf{Model B:} Risk features $+$ public controls.
\end{enumerate}

Compare the information coefficient ($\IC$) and SHAP importance
(Chapter~\ref{ch:interpretation}) of each risk feature across
Models A and B.

\begin{keyidea}[Confound Checking Is Not Optional]
If your signal vanishes after adding public controls, it is redundant
with publicly available information---you have not added value.  The
test is simple: compare $\IC$ with and without controls.  A feature
that survives controls has \emph{incremental} information from
proprietary risk data.  A feature that does not is just noise around
VIX or credit spreads.

This is the feature-level analogue of the confound checks in
asset pricing (Chapter~\ref{ch:asset-pricing}): a ``new'' factor that
is spanned by existing factors adds nothing.
\end{keyidea}


% ──────────────────────────────────────────────────────────────
\section{Summary}
\label{sec:features-summary}
% ──────────────────────────────────────────────────────────────

Table~\ref{tab:feature-families} maps each feature family to its
constituent features, source chapter, and project phase.

\begin{table}[ht]
\centering
\small
\begin{tabular}{@{}p{2.6cm}p{4.5cm}p{2.8cm}p{2.8cm}@{}}
\toprule
\textbf{Feature Family} & \textbf{Features} &
  \textbf{Source Chapter} & \textbf{Project Phase} \\
\midrule
VaR Utilization
  & Level, $\Delta h$-day change, z-score
  & Ch.~\ref{ch:risk-systems} \S\ref{sec:var-utilization}
  & Phase 2 (priority) \\[6pt]

Factor Concentration
  & HHI, top-3 share, $\Delta$HHI
  & Ch.~\ref{ch:risk-systems} \S\ref{sec:factor-var}
  & Phase 2 (priority) \\[6pt]

VaR Dynamics
  & $\Delta$VaR, CVaR by asset class, VaR momentum
  & Ch.~\ref{ch:risk-systems} \S\ref{sec:var-methods}
  & Phase 2 \\[6pt]

Scenario P\&L
  & Worst-case id, dispersion, flip indicator, skewness
  & Ch.~\ref{ch:risk-systems} \S\ref{sec:scenarios}
  & Phase 2 \\[6pt]

Cross-Asset Flow
  & CVaR share shifts, rolling cross-asset correlation
  & Ch.~\ref{ch:risk-systems} \S\ref{sec:component-var}
  & Phase 2 \\[6pt]

Dealer Greeks
  & Net $\Gamma$/$\mathcal{V}$/vanna/charm, flip level,
    distance to flip
  & Ch.~\ref{ch:microstructure} \S\ref{sec:gamma}
  & Phase 4B \\
\bottomrule
\end{tabular}
\caption{Feature family reference.  Priority families (VaR utilization
  and factor concentration) are tested first in Phase~2 per the design
  spec.  All batch-computed features require PIT lagging
  (Section~\ref{sec:pit}).  Dealer Greeks are real-time and belong to
  Project~2.}
\label{tab:feature-families}
\end{table}

\vspace{1em}

\noindent\textbf{Key takeaways:}
\begin{enumerate}[nosep]
  \item Point-in-time discipline is non-negotiable.  Enforce PIT lagging
        at the data ingestion layer.  VaR for date $T$ is available at
        $T+1$ open.
  \item Six feature families span the full set of SecDB risk outputs.
        VaR utilization and factor concentration are tested first
        (strongest theory: Adrian and Shin, 2014; He, Kelly, and
        Manela, 2017; Coval and Stafford, 2007).
  \item Every feature must pass confound checks against public controls
        (VIX, credit spread, term slope).  Features that do not survive
        controls are noisy proxies, not new signals.
  \item The level, rate of change, and z-score of a quantity are three
        different features---they capture different aspects
        (constraint state, direction, unusualness) and should all be
        included.
  \item Dealer Greek features (Project~2) are real-time and do not
        require PIT lagging, but they do require a consistent intraday
        timestamp for reproducibility.
\end{enumerate}
